\documentclass[12pt, twoside]{article}
%\documentclass[12pt, twoside]{article}
\usepackage[letterpaper, margin=1in, headsep=0.2in]{geometry}
\setlength{\headheight}{0.6in}
%\usepackage[english]{babel}
\usepackage[utf8]{inputenc}
\usepackage{microtype}
\usepackage{amsmath}
\usepackage{amssymb}
%\usepackage{amsfonts}
\usepackage[nomessages]{fp} %\FPeval{\var-name}{2*sin(pi/6)}
\usepackage{siunitx} %units in math. eg 20\milli\meter
\usepackage{yhmath} % for arcs, overparenth command
\usepackage{tikz} %graphics
\usetikzlibrary{quotes, angles, arrows, arrows.meta}
\usepackage{pgfplots}
\usepgfplotslibrary{fillbetween}
\pgfplotsset{compat=1.16}
\usepackage{graphicx} %consider setting \graphicspath{{images/}}
\usepackage{parskip} %no paragraph indent
\usepackage{enumitem}
\usepackage{multicol}
\usepackage{venndiagram}

\usepackage{fancyhdr}
\pagestyle{fancy}
\fancyhf{}
\renewcommand{\headrulewidth}{0pt} % disable the underline of the header
\raggedbottom
\hfuzz=2mm %suppresses overfull box warnings

\usepackage{hyperref}

\fancyhead[LE]{\thepage}
\fancyhead[RO]{Markscheme}
\fancyhead[LO]{La Scuola d'Italia / Dr.\ Huson / IB Math 11 \\* 3 November 2026}

\begin{document}
\subsubsection*{2.8 Classwork: Systems of Equations by Elimination --- Markscheme}

\begin{enumerate}[itemsep=0.3cm]

%--- Problem 1: Three 2x2 systems by elimination: add directly; scale one equation; scale both equations [9 marks] ---
\item[1.] [Maximum mark: 9]

\emph{In each part: M1 for adding or subtracting the equations (after
scaling, if used) so that one variable is eliminated, seen as the
resulting one-variable equation. Substitution in place of elimination
earns M0, then A1A1 for correct values. Correct values with no working
(for example from the GDC) earn M0A1A1. FT the second value from the
candidate's first value.}

\textbf{(a)} $3x + 2y = 17$, \ $5x - 2y = 7$

Adding: $8x = 24$ \hfill \textbf{M1}

$x = 3$ \hfill \textbf{A1}

$3(3) + 2y = 17$, so $y = 4$ \hfill \textbf{A1}

\textbf{(b)} $5x + 3y = 21$, \ $4x - y = 27$

$3 \times$ second equation: $12x - 3y = 81$

Adding: $17x = 102$ \hfill \textbf{M1}

$x = 6$ \hfill \textbf{A1}

$5(6) + 3y = 21$, so $3y = -9$ and $y = -3$ \hfill \textbf{A1}

\emph{Common error: scaling the left side only, $12x - 3y = 27$, gives
$17x = 48$; M1 is not earned because the scaled equation is wrong.
Accept eliminating $x$ first: $4(5x + 3y) - 5(4x - y)$ gives
$17y = -51$.}

\textbf{(c)} $3x + 4y = 14$, \ $2x + 3y = 11$

$2 \times$ first: $6x + 8y = 28$; \quad $3 \times$ second: $6x + 9y = 33$

Subtracting: $(6x + 9y) - (6x + 8y) = 33 - 28$ \hfill \textbf{M1}

$y = 5$ \hfill \textbf{A1}

$3x + 4(5) = 14$, so $3x = -6$ and $x = -2$ \hfill \textbf{A1}

\emph{Accept eliminating $y$ first: $3(3x + 4y) - 4(2x + 3y)$ gives
$x = 42 - 44 = -2$. Common error: $x = 2$ (sign slip in $14 - 20$) earns
the final A0.}

\newpage

%--- Problem 2: Supply and demand lines from survey data; show demand equation; find supply equation; equilibrium by elimination; minimum price surplus [12 marks] ---
\item[2.] [Maximum mark: 12]

Survey points $(q,\ p)$: demand $(16,\ 36)$ and $(4,\ 72)$; supply
$(6,\ 33)$ and $(20,\ 68)$. Quantity $q$ in hundreds of helmets per week,
price $p$ in euros.

\textbf{(a)} $m = \dfrac{72 - 36}{4 - 16} = \dfrac{36}{-12} = -3$
\hfill \textbf{M1}

$36 = -3(16) + c$, so $c = 84$ \hfill \textbf{A1}

$p = 84 - 3q$ \hfill \textbf{AG}

\emph{M1 for the gradient from the two demand points. A1 for $c = 84$
from a point. Only checking that the given line fits both points, with
no gradient, earns M1A0.
A gradient of $-\frac{1}{3}$ (run over rise) earns M0A0.}

\textbf{(b)} $m = \dfrac{68 - 33}{20 - 6} = \dfrac{35}{14}$ \hfill \textbf{M1}

$= 2.5$ \hfill \textbf{A1}

$33 = 2.5(6) + c$, so $c = 18$: \quad $p = 2.5q + 18$ \hfill \textbf{A1}

\emph{Accept $p - 33 = 2.5(q - 6)$ or any equivalent form. FT the
equation from the candidate's gradient. Wrong method: the gradient
inverted, $\frac{14}{35} = 0.4$, earns M0A0, then A1 FT for
$p = 0.4q + 30.6$.}

\textbf{(c)} $3q + p = 84$ and $2.5q - p = -18$ \hfill \textbf{A1}

Adding: $5.5q = 66$ \hfill \textbf{M1}

$q = 12$, so $1200$ helmets per week \hfill \textbf{A1}

$p = 84 - 3(12) = 48$ (euros) \hfill \textbf{A1}

\emph{A1 for both equations in the form $aq + bp = c$; accept
$5q - 2p = -36$ or any equivalent pair. M1 for adding or subtracting
(after scaling, if used) so that one variable is eliminated. Accept
eliminating $q$ first: $5(3q + p) - 3(5q - 2p)$ gives $11p = 528$.
Substitution in place of elimination earns M0, then A1A1. Accept $q = 12$
(hundred) for the first value. FT from the candidate's supply line
in (b).}

\emph{Wrong methods: a sign slip, $2.5q - p = 18$, gives $5.5q = 102$,
$q = 18.5$ (18.5454\ldots), earning A0M1A0A0.}

\textbf{(d)} Substituting $p = 60$ into both equations \hfill \textbf{M1}

Demand: $60 = 84 - 3q$, so $q = 8$. Supply: $60 = 2.5q + 18$, so
$q = 16.8$. \hfill \textbf{A1}

Surplus $= 16.8 - 8 = 8.8$ hundred, so $880$ helmets per week
\hfill \textbf{A1}

\emph{The first A1 needs both quantities. Accept $8.8$ (hundred) or $880$
for the surplus. A negative answer, or calling it a shortage, earns the
final A0. FT the surplus from the candidate's supply line in (b), provided
supply exceeds demand. Wrong methods: comparing with the equilibrium
quantity, $16.8 - 12 = 4.8$ or $12 - 8 = 4$, earns the final A0.}

\end{enumerate}
\end{document}
